% ---- pg-564 （书页 552） ----

$q\le N-1$ 维空间中的某些结构, 使上述匹配尽可能良好. 匹配优度的数值量度称为应力.

有可能将 $N$ 个项目在一个低维坐标系中作出分类, 且只用 $N(N-1)/2$ 个原始相似性(距离)的\mvbold{等级顺序}而不用其数值. 当只用这种顺序信息来获得几何表示时, 这个过程称为\mvbold{非度量性多维标度变换}. 如果利用原始相似性(距离)的实际数值来获得一个 $q$ 维的几何表示, 则此种过程称为\mvbold{度量性多维标度变换}. 度量性多维标度变换亦称\mvbold{主坐标分析}.

标度变换方法为谢泼德(Shepard, 见文献[29]中对早期工作的评论)和克鲁斯卡尔(Kruskal, 见文献[19,20,21])等人所提出. 关于多维标度变换的历史、理论和应用, 文献[35]提供了一个很好的总结. 多维标度变换法需要使用计算机. 目前有几个好的计算机程序可用于这一目的.

\mvsub{基本算法}

对 $N$ 个项目, 各对不同项目之间共有 $M=N(N-1)/2$ 个相似性(距离). 这些相似性构成基本数据(如果相似性不易量化, 例如两种颜色间的相似性, 则相似性的等级顺序便是基本数据).

假定没有结(相等), 则相似性数据可依严格上升的顺序排成一列:
\phantomsection\label{eq:12-21}
\begin{equation}
s_{i_1k_1} < s_{i_2k_2} < \cdots < s_{i_Mk_M} \tag{12-21}
\end{equation}
这里 $s_{i_1k_1}$ 为 $M$ 个相似性中最小者. 下标 $i_1k_1$ 指示最不相似的那一对项目(即在相似性排序中等级为 1 的项目). 其他下标可按类似方式解释. 我们要找到这 $N$ 个项目的一个 $q$ 维点结构, 使得各对项目之间的距离 $d_{ik}^{(q)}$ 与\hyperref[eq:12-21]{式(12-21)}中的排序相匹配. 当距离 $d_{ik}^{(q)}$ 的排序与\hyperref[eq:12-21]{式(12-21)}中的排序相对应时, 或者说当
\phantomsection\label{eq:12-22}
\begin{equation}
d_{i_1k_1}^{(q)} > d_{i_2k_2}^{(q)} > \cdots > d_{i_Mk_M}^{(q)} \tag{12-22}
\end{equation}
时, 便得到一个完美的匹配. 就是说, $q$ 维空间中距离的递降顺序与原始相似性的递增顺序完全类似. 这里重要的只是\hyperref[eq:12-22]{式(12-22)}中的排序, 距离的大小是无关紧要的.

对给定的 $q$ 值, 可能找不到这样一种 $q$ 维空间的结构, 其点间距离与原始相似性之间有单调关系. 克鲁斯卡尔(见文献[19])曾提出一个用以衡量某个几何表示偏离完美匹配的程度的量度, 这个量度称为应力. 其定义为
\phantomsection\label{eq:12-23}
\begin{equation}
\text{应力}(q) = \left\{ \frac{\sum\limits_{i<k}\sum (d_{ik}^{(q)} - \hat{d}_{ik}^{(q)})^2}{\sum\limits_{i<k}\sum [d_{ik}^{(q)}]^2} \right\}^{1/2} \tag{12-23}
\end{equation}
其中 $\hat{d}_{ik}^{(q)}$ 满足\hyperref[eq:12-22]{式(12-22)}, 即它们与相似性之间有单调关系. 这些 $\hat{d}_{ik}^{(q)}$ 并不是真正的距离, 因为它们不满足\hyperref[eq:1-25]{式(1-25)}中距离的常见性质. 它们仅仅是用来评价真实距离 $d_{ik}^{(q)}$ 的非单调性的一些参考数值.

基本的想法是找到 $q$ 维空间中 $N$ 个项目的点表示, 使应力$(q)$尽可能地小. 克鲁斯卡尔(见文献[19])建议根据以下准则对应力作非正式解释:

\begin{mvtable}
\begin{tabular}{cc}
\toprule
应力 & 拟合优度\\
\midrule
20\% & 差\\
10\% & 一般\\
5\% & 好\\
2.5\% & 非常好\\
0\% & 完美\\
\bottomrule
\end{tabular}
\end{mvtable}

\noindent\hfill (12-24)\hspace*{1em}
